\chapter{Source code and data analysis}
\label{cha:appendix_code}

The program code developed for this thesis is available through the following \texttt{GitHub} repositories.
It supports data analysis, data fitting,
derived physical-property calculations, and figure generation.

\hyperref[cha:project_bilayers]{Chapter~\ref*{cha:project_bilayers}: \nameref*{cha:project_bilayers}}
\cite{niu2024electronic}

\quad\href{https://github.com/Photonico/Graphene-BC_20230622}
{\texttt{https://github.com/Photonico/Graphene-BC\_20230622}}

\hyperref[cha:project_boron]{Chapter~\ref*{cha:project_boron}: \nameref*{cha:project_boron}}
\cite{niu2026superconductivity}

\quad\href{https://github.com/Photonico/o-B14_20241024}
{\texttt{https://github.com/Photonico/o-B14\_20241024}}

\hyperref[cha:project_beryllene]{Chapter~\ref*{cha:project_beryllene}: \nameref*{cha:project_beryllene}}
\cite{niu2026draftemergent}

\quad\href{https://github.com/Photonico/H-Beryllene_20250718}
{\texttt{https://github.com/Photonico/H-Beryllene\_20250718}}

The thesis source files and related materials:

\quad\href{https://github.com/Photonico/PhD_thesis_20251216}
{\texttt{https://github.com/Photonico/PhD\_thesis\_20251216}}

Each chapter repository contains the scripts and \texttt{Jupyter} notebooks used for the corresponding chapter.
The main reusable analysis code is stored in a top-level \texttt{vmatplot} directory within each project repository.
These local versions share the same code base but reflect different stages of development.
Later repositories contain newer functionality and refinements.
The \texttt{Jupyter} notebooks provide complete usage examples for the corresponding analyses.

The candidate used \texttt{vmatplot} to process first-principles calculation data,
analyse band structures, total and projected densities of states, dielectric functions, phonon dispersions, AIMD trajectories, and convergence tests, fit numerical data, and generate figures.
In particular, the code calculates the derived optical properties from the dielectric function.

Unless otherwise stated in the repository,
the candidate releases the source code and related scripts under the MIT License.
This permissive licence allows reuse, modification, distribution, and adaptation of the code,
with retention of the original copyright and licence notice.
